\subsubsection{Scope and sufficiency}
The selected HumanEval Input 1 response gives B2 a rating of 10 while saying that a checking procedure for HumanEval answers is not reported. It cites verification for a separate training set. The Input 2 response instead uses insufficient evidence. This is an example of why an excerpt must concern the benchmark and check being judged; locating the words in the paper is not enough.

For HumanEval B4, the Input 1 response uses not applicable, while the Input 2 response uses 5 because the paper discusses expected student performance. Both received the same paper. The difference concerns whether the review includes that wider comparison or only measured benchmark results. The present reference proposal does not independently settle that scope choice.

Other changes also concern thresholds applied to shared text. For DesignQA A5, the Input 2 response lowers the rating while asking for model inference steps and expected outputs for every task. For TruthfulQA A3, the Input 2 response asks for a fuller partition scheme. These changes are not automatically corrections: the relevance of the added requirements needs explanation. Conversely, the selected TruthfulQA B2 responses both document validation disagreements without treating the existence of disagreements as absence of checking.

These observations cover all sixty judgments in the six preselected responses through an accompanying audit. They are interpretations by the preparing assistant, not independent human error counts. No semantic accuracy rate is inferred from quotation matching or from the selected examples. The Claude batch is reported separately below.
